% Źródła: Eves s. 283--284; Greenberg 2010, s. 205, 207, 210; Wikipedia: Parallel postulate, Proclus, Clavius, Wallis.
% https://en.wikipedia.org/wiki/Parallel_postulate

\section{Piąty postulat Euklidesa} % lang-pl
\section{Il quinto postulato di Euclide} % lang-it
\label{sekcja_piaty_postulat}

W tej sekcji zobaczymy, dlaczego jeden z pięciu postulatów Euklidesa będzie budził podejrzenia przez ponad dwa tysiące lat i kto będzie próbował się go pozbyć. % lang-pl
In questa sezione vedremo perché uno dei cinque postulati di Euclide susciterà sospetti per più di duemila anni e chi cercherà di sbarazzarsene. % lang-it

\begin{axiom}
[piąty postulat Euklidesa] % lang-pl
[quinto postulato di Euclide] % lang-it
\index{postulat!piąty (równoległości)} % lang-pl
\index{postulato!quinto (delle parallele)} % lang-it
    Jeżeli prosta przecinająca dwie proste tworzy z nimi po jednej stronie kąty wewnętrzne o sumie mniejszej niż dwa kąty proste, to te dwie proste, przedłużone nieograniczenie, przetną się po tej stronie, po której suma kątów jest mniejsza niż dwa kąty proste. % lang-pl
    Se una retta che ne interseca altre due forma con esse, da una stessa parte, angoli interni la cui somma è minore di due angoli retti, allora le due rette, prolungate indefinitamente, si incontrano da quella parte in cui la somma degli angoli è minore di due angoli retti. % lang-it
\end{axiom}

Eves \cite[s. 283]{eves1_1972}. % lang-pl
Eves \cite[p. 283]{eves1_1972}. % lang-it

Wystarczy porównać go z czterema pierwszymi postulatami, żeby zobaczyć różnicę: brakuje mu zwięzłości i oczywistości, a w gruncie rzeczy przypomina odwrócenie twierdzenia I.17. % lang-pl
Euklides sam wyraźnie będzie go unikał i użyje go dopiero w dowodzie twierdzenia I.29. % lang-pl
Nie dziwi więc, że wielu greckich geometrów uzna go raczej za twierdzenie niż za postulat i zapragnie wyprowadzić go z pozostałych albo zastąpić czymś bardziej przekonującym \cite[s. 283--284]{eves1_1972}. % lang-pl
Basta confrontarlo con i primi quattro postulati per vedere la differenza: gli manca la concisione e l'evidenza, e in fondo ricorda l'inverso della proposizione I.17. % lang-it
Lo stesso Euclide lo eviterà con evidente cura e lo userà solo nella dimostrazione della proposizione I.29. % lang-it
Non stupisce dunque che molti geometri greci lo giudichino più un teorema che un postulato e desiderino dedurlo dagli altri o sostituirlo con qualcosa di più convincente \cite[p. 283--284]{eves1_1972}. % lang-it

Ptolemeusz spróbuje udowodnić postulat, a jego rozumowanie Proklos uzna za nieudane. % lang-pl
Sam Proklos (410--485) oprze swój dowód na zasadzie Arystotelesa: odcinki prostopadłe opuszczone z jednego ramienia kąta ostrego na drugie nie mają ograniczenia z góry. % lang-pl
Zasada ta okaże się jednak prawdziwa także w geometrii hiperbolicznej, więc sama nie wystarczy do dowodu -- Greenberg \cite[s. 205, 207, 210]{greenberg_2010} zauważa, że Proklos pierwszy dostrzeże jej znaczenie dla podstaw geometrii. % lang-pl
Jego komentarz do pierwszej księgi \emph{Elementów} stanie się jednym z najcenniejszych źródeł do historii starożytnej matematyki, bo zachowa też wiadomości o Geminosie i Posejdoniosie. % lang-pl
Tolomeo tenterà di dimostrare il postulato, e il suo ragionamento sarà giudicato da Proclo insufficiente. % lang-it
Lo stesso Proclo (410--485) fonderà la sua dimostrazione sul principio di Aristotele: i segmenti perpendicolari condotti da un lato di un angolo acuto all'altro non hanno limite superiore. % lang-it
Questo principio si rivelerà però vero anche nella geometria iperbolica, e dunque da solo non basta per la dimostrazione -- Greenberg \cite[p. 205, 207, 210]{greenberg_2010} osserva che Proclo sarà il primo a riconoscerne l'importanza per i fondamenti della geometria. % lang-it
Il suo commento al primo libro degli \emph{Elementi} diventerà una delle fonti più preziose per la storia della matematica antica, perché conserverà anche notizie su Gemino e Posidonio. % lang-it
\index[persons]{Proklos} % lang-pl
\index[persons]{Proclo} % lang-it

Pałeczkę przejmą uczeni świata islamu. % lang-pl
Ibn al-Hajsam (965--1039) poprowadzi dowód nie wprost, wprowadzając do geometrii ruch, i rozważy czworokąt z trzema kątami prostymi. % lang-pl
Omar Chajjam (ok. 1050--1123) oprze swój dowód na zasadzie Arystotelesa: dwie zbieżne proste przecinają się. % lang-pl
Nasir ad-Din at-Tusi (1201--1274) napisze w 1250 roku traktat \emph{Dyskusja usuwająca wątpliwości co do linii równoległych}, a w 1298 roku jego syn Sadr ad-Din przedstawi hipotezę równoważną postulatowi. % lang-pl
Il testimone passerà agli studiosi del mondo islamico. % lang-it
Ibn al-Haytham (965--1039) condurrà una dimostrazione per assurdo, introducendo in geometria il movimento, e considererà un quadrilatero con tre angoli retti. % lang-it
Omar Khayyam (1050 circa--1123) fonderà la sua dimostrazione sul principio di Aristotele: due rette convergenti si intersecano. % lang-it
Nasir al-Din al-Tusi (1201--1274) scriverà nel 1250 il trattato \emph{Discussione che rimuove il dubbio sulle rette parallele}, e nel 1298 suo figlio Sadr al-Din presenterà un'ipotesi equivalente al postulato. % lang-it
\index[persons]{Ibn al-Hajsam} % lang-pl
\index[persons]{Ibn al-Haytham} % lang-it
\index[persons]{Omar Chajjam} % lang-pl
\index[persons]{Omar Khayyam} % lang-it
\index[persons]{Nasir ad-Din at-Tusi} % lang-pl
\index[persons]{Nasir al-Din al-Tusi} % lang-it

W Europie do sprawy wrócą Christopher Clavius, jezuita z Collegio Romano, w swoim wydaniu Euklidesa z 1574 roku, Giordano Vitale w książce \emph{Euclide restituto} (1680), a John Wallis w wykładzie z 1663 roku, do którego zainspiruje go traktat Sadr ad-Dina. % lang-pl
Każdy z nich zostawi po sobie zdanie, które zastępuje piąty postulat -- o nich opowiemy w następnej sekcji. % lang-pl
In Europa riprenderanno la questione Christopher Clavius, gesuita del Collegio Romano, nella sua edizione di Euclide del 1574, Giordano Vitale nel libro \emph{Euclide restituto} (1680) e John Wallis in una lezione del 1663, ispirata dal trattato di Sadr al-Din. % lang-it
Ciascuno di loro lascerà un enunciato che sostituisce il quinto postulato -- ne parleremo nella sezione seguente. % lang-it
\index[persons]{Clavius, Christopher} % lang-pl
\index[persons]{Clavio, Cristoforo} % lang-it
\index[persons]{Vitale, Giordano}%
\index[persons]{Wallis, John}%

% Wikipedia podaje daty życia Chajjama różnie (XI/XII w.); do sprawdzenia w źródle drukowanym.
